\documentclass[11pt]{article}
\usepackage[margin=1in]{geometry}
\usepackage{amsmath,amssymb,amsthm}
\usepackage[hidelinks]{hyperref}
\newtheorem{theorem}{Theorem}
\newtheorem{lemma}{Lemma}
\title{Positivity of the Fibonacci characteristic function\\
\large A solution to MathDB 360139}
\author{Principia Math In Progress}
\date{18 September 2026}
\begin{document}
\maketitle
\begin{abstract}
The characteristic generating function of the letter $a$ in the
Fibonacci fixed word is strictly positive on $(-1,1)$. We prove this
by telescoping periodic-prefix approximants and bounding their negative
increments using the spacing of Fibonacci numbers. An elementary
instance of the source's even-block decomposition handles the remaining
compact interval. This proves Conjecture 5.12 of Pouti, Ramsey, and
Strungaru.
\end{abstract}

\section*{Task 1: the original problem and its full solution}
Let $\sigma(a)=ab$, $\sigma(b)=a$, and let
$w=w_0w_1\cdots=abaababa\cdots$ be its fixed word starting in $a$.
The characteristic function in \cite[Sections 3--5]{PRS} is
\[
 C_a(z)=\sum_{j\geq0}{\bf1}_{\{w_j=a\}}z^j,\qquad |z|<1.
\]
In particular the constant coefficient is $1$. All series below
are considered inside their disk of absolute convergence.

\begin{theorem}\label{thm:main}
For every real $x\in(-1,1)$, one has $C_a(x)>0$.
Moreover, for $9/10\leq t<1$,
\[
 C_a(-t)>\frac{1}{1+t}-\frac{127(1+t)}{600t^2}
 \geq\frac{2753}{48600(1+t)}>\frac{2753}{97200}.
\]
\end{theorem}

\subsection*{An exact telescoping expansion}
Write $A_{-1}=b$, $A_0=a$, and $A_{n+1}=A_nA_{n-1}$ for $n\geq0$.
Then $A_n=\sigma^n(a)$ for $n\geq0$. Put
$L_n=|A_n|=F_{n+2}$, where $F_0=0,F_1=1$, and set
\[
 P_n(z)=\sum_{\substack{0\leq j<L_n\\(A_n)_j=a}}z^j,
 \qquad f_n(z)=\frac{P_n(z)}{1-z^{L_n}}\quad(n\geq0).
\]
Here $P_{-1}=0$, $P_0=1$, and
\begin{equation}\label{eq:rec}
 P_{n+1}=P_n+z^{L_n}P_{n-1}.
\end{equation}
Since the $A_n$ are increasing prefixes of $w$,
$P_n(z)\to C_a(z)$ and $f_n(z)\to C_a(z)$ for $|z|<1$.

Define
\[
 D_n=P_n(1-z^{L_{n-1}})-P_{n-1}(1-z^{L_n}),\qquad n\geq0.
\]
Equation \eqref{eq:rec} gives $D_0=1-z$ and
$D_{n+1}=-z^{L_n}D_n$. Since
$\sum_{j=0}^{n-1}L_j=L_{n+1}-2$, induction proves
\begin{equation}\label{eq:D}
 D_n=(-1)^nz^{L_{n+1}-2}(1-z).
\end{equation}
Let $a=L_n$, $b=L_{n-1}$, $c=2a+b=L_{n+2}$, and
$d=3a+2b=L_{n+3}=2c-a$. The word identity
$A_{n+3}=A_nA_{n-1}A_nA_nA_{n-1}$ gives
\[
 P_{n+3}=P_n(1+z^{a+b}+z^{2a+b})
             +P_{n-1}(z^a+z^{3a+b}).
\]
Subtracting after a common denominator and using \eqref{eq:D} yields
\begin{align}
 P_{n+3}(1-z^a)-P_n(1-z^d)
 &=-z^a(1+z^c)D_n \notag\\
 &=(-1)^{n+1}z^{c-2}(1+z^c)(1-z).
 \label{eq:triple}
\end{align}

For $n=3k$, all of $a,b,c,d$ are odd, because $F_m$ is even
exactly when $3$ divides $m$. Thus, for $0<t<1$,
\[
 f_{3k+3}(-t)-f_{3k}(-t)=(-1)^k E_k(t),
\]
where
\begin{equation}\label{eq:E}
 E_k(t)=
 \frac{(1+t)t^{c_k-2}(1-t^{c_k})}
 {(1+t^{a_k})(1+t^{2c_k-a_k})}>0,\qquad
 a_k=F_{3k+2},\quad c_k=F_{3k+4}.
\end{equation}
The bound $E_k(t)\leq2t^{c_k-2}$ makes the series of increments
absolutely convergent for each fixed $t<1$: the strictly increasing
integers $c_k$ satisfy $c_k\geq k+3$. Taking limits in the telescoping
sum therefore gives the exact identity
\begin{equation}\label{eq:series}
 C_a(-t)=\frac{1}{1+t}+\sum_{k\geq0}(-1)^kE_k(t).
\end{equation}

\subsection*{A bound for widely spaced negative increments}
\begin{lemma}\label{lem:spacing}
Let $x_0,x_1,\ldots>0$ satisfy $x_{\ell+1}\geq16x_\ell$, and put
$g(u)=u(1-u)/(1+u)^2$. Then
\[
 \sum_{\ell\geq0}g(e^{-x_\ell})<\frac{127}{600}.
\]
\end{lemma}
\begin{proof}
For $0<u<1$, the inequality $g(u)\leq1/8$ follows from
\[
 \frac18-g(u)=\frac{(1-3u)^2}{8(1+u)^2}\geq0.
\]
For $x>0$, also
\[
 g(e^{-x})=
 \frac{e^{-x}}{1+e^{-x}}\tanh(x/2)\leq\frac{x}{4},
 \qquad
 g(e^{-x})\leq e^{-x}.
\]
Here $\tanh y\leq y$ for $y\geq0$, as follows by differentiating.
Split the terms into $x<1/4$, $1/4\leq x<4$, and $x\geq4$.
The first group is finite, and summing backwards from its largest
element gives
\[
 \sum_{x_\ell<1/4}g(e^{-x_\ell})
 <\frac14\,\frac{1/4}{1-1/16}=\frac1{15}.
\]
The same strict bound holds if this group is empty.
The middle group contains at most one element, hence contributes
at most $1/8$. In the last group, the successive arguments are at
least $4\cdot16^r\geq4+60r$, so its contribution is at most
\[
 \sum_{r\geq0}e^{-4-60r}=\frac{e^{-4}}{1-e^{-60}}<\frac1{50}.
\]
Indeed, $e^4>\sum_{j=0}^7 4^j/j!=16319/315>51$ and $e^{60}>51$,
which imply the last strict inequality. Adding the three bounds gives
$1/15+1/8+1/50=127/600$.
\end{proof}

By the arithmetic--geometric mean inequality,
\[
 (1+t^a)(1+t^{2c-a})\geq(1+t^c)^2.
\]
Consequently \eqref{eq:E} implies
\begin{equation}\label{eq:majorant}
 E_k(t)\leq\frac{1+t}{t^2}g(t^{c_k}).
\end{equation}
The negative increments in \eqref{eq:series} have
$c_{2\ell+1}=F_{6\ell+7}$. For $m\geq7$,
\[
 F_{m+6}=13F_m+8F_{m-1}>17F_m>16F_m,
\]
since $F_{m-1}>F_m/2$. Thus
$x_\ell=-F_{6\ell+7}\log t$ meets Lemma \ref{lem:spacing}.
Discarding the positive increments in \eqref{eq:series} gives
\begin{equation}\label{eq:lower}
 C_a(-t)>
 \frac1{1+t}-\frac{127(1+t)}{600t^2}.
\end{equation}
For $9/10\leq t<1$, $(1+t)/t\leq19/9$ and
\[
 1-\frac{127}{600}\left(\frac{19}{9}\right)^2
 =\frac{2753}{48600}>0.
\]
This proves all the quantitative assertions of Theorem \ref{thm:main}.

\subsection*{The remaining compact interval}
We now reprove the needed low-order case of the even-block
decomposition in \cite[Section 5]{PRS}; that decomposition is an
established input, not a new claim.
There is no occurrence of $bb$ in $w$: every substitution image starts
in $a$, and a $b$ only occurs at the end of $\sigma(a)$.
Since $w=\sigma^3(w)$, group the letters of the latter occurrence of
$w$ into successive pairs, starting at its beginning, and apply
$\sigma^3$. With $A=\sigma^3(a)=abaab$ and $B=\sigma^3(b)=aba$,
this expresses $w$ as a concatenation of the blocks $AB,AA,BA$.
Each block has even length. Their $a$-characteristic polynomials
evaluated at $-t$ are, respectively,
\begin{align*}
 R(t)&=1+t^2-t^3-t^5-t^7,\\
 S(t)&=(1+t^2-t^3)(1-t^5),\\
 T(t)&=(1+t^2)(1-t^3)+t^6.
\end{align*}
The latter two are positive for $0\leq t<1$. Also
\[
 R'(t)=t(2-3t-5t^3-7t^5).
\]
The parenthesis is strictly decreasing on $[0,\infty)$; thus $R$
first increases and then decreases, and its minimum on $[0,9/10]$
occurs at an endpoint. Both endpoint values are positive:
\[
 R(0)=1,\qquad R(9/10)=\frac{122131}{10000000}.
\]
All offsets of the concatenated blocks are even. Hence their
contributions to $C_a(-t)$ are nonnegative powers of $t$ times
positive numbers for $0\leq t\leq9/10$. The first block already
contributes a positive number. Absolute convergence permits this
grouping, proving $C_a(-t)>0$ on this interval.

Finally, $C_a(x)\geq1$ for $0\leq x<1$, because its coefficients
are nonnegative and its constant term is $1$.
Together with \eqref{eq:lower}, this proves Theorem \ref{thm:main}
on the entire required interval.

\section*{Tasks 2--5: proposed follow-up directions}
These are questions for further work, not additional solved claims.
\begin{description}
\item[Task 2.] Determine the set of limit points of $C_a(-t)$ as
$t\to1^-$, including exact or rigorously characterized values of
its lower and upper limits.
\item[Task 3.] Determine the optimal constant
$c_*=\inf_{0\leq t<1}C_a(-t)$ and whether that infimum is attained
at an interior point.
\item[Task 4.] For each integer $m\geq2$, let $w^{(m)}$ be the fixed
word starting in $a$ of $a\mapsto a^mb$, $b\mapsto a$.
Classify the $m$ for which its $a$-characteristic generating
function is positive throughout $(-1,1)$.
\item[Task 5.] Determine the largest $\rho\in(0,1]$ for which
$C_a(z)$ has no zero in $|z|<\rho$, and, if $\rho<1$, rigorously
locate and count all zeros with modulus $\rho$.
\end{description}

\paragraph{Verification and status.}
The accompanying verifier checks the polynomial identities against
directly substituted words, the exact constants and block
polynomials, and finite telescoping sums in rational arithmetic.
These checks supplement the analytic proof for the infinite series.
The current 23-page source was read in full; its authors list the
paper as accepted or in press in 2026. The dated literature audit
records the bounded prior-claim search. Verification is by the
authoring agent only, without independent peer review.

\begin{thebibliography}{9}
\raggedright
\bibitem{PRS} A. Pouti, C. Ramsey, and N. Strungaru,
\emph{Generating functions of substitutions},
\href{https://arxiv.org/abs/2309.16100}{arXiv:2309.16100v2},
25 May 2026, 23 pp.; accepted/in press, Rocky Mountain Journal of
Mathematics, according to the authors' publication lists checked
18 September 2026. The target is Conjecture 5.12 on p.~21.
\end{thebibliography}
\end{document}
